% claim-ref: cherry-shrimp#published-ph
% claim-ref: cherry-shrimp#hardness-minima
% claim-ref: cherry-shrimp#specialist-chemistry
% claim-ref: cherry-shrimp#test-hardness-separately
% claim-ref: cherry-shrimp#nursery-proposal
% claim-ref: cherry-shrimp#water-gate
% claim-ref: cherry-shrimp#substrate-intake
% claim-ref: cherry-shrimp#plant-plan
% claim-ref: cherry-shrimp#possible-temperature-overlap
% claim-ref: cherry-shrimp#nursery-cycle-status
\CompanionHeader{TANK SETUP}{Cherry shrimp}{Neocaridina davidi / measured water before changes}{The possible 10-gallon nursery is not yet verified cycled or ready.}
\Context{Published guidance: Co-Op pH 6.5--8.5, GH at least 6 dGH and KH at least 2 dKH. These are a broad range and minima, not one water recipe or household measurements. \textbf{[CS-COOP]}}
\CardRow{Keep advice separate}{The Shrimp Farm's narrower ideals: pH 7.0--7.6, GH 9--11 and KH 4--6 degrees. Check supplier and replacement water; do not abruptly retune a healthy aquarium to a table. \textbf{[CS-FARM]}}{GH and KH differ}{GH mainly reflects calcium and magnesium; KH reflects carbonate/bicarbonate buffering. Measure both separately in degrees. A pH reading cannot replace either test. \textbf{[CS-GH; CS-KH]}}
\CardRow{Cycle without animals}{Proposed gate: establish and document biofiltration without live shrimp. Require undetectable ammonia/nitrite on working tests; log nitrate and units. One clean test or elapsed days is not cycle proof. \textbf{[CS-VET]}}{Substrate and intakes}{Inert substrate can suit this species. Use a sponge filter or protected intake and maintain flow. Review gaps and suction for tiny juveniles; an air stone alone proves neither filtration nor oxygen adequacy. \textbf{[CS-FARM; proposed]}}
\Panel{Plants and grazing surfaces / proposed}{Healthy hornwort and guppy grass can provide cover and grazing surfaces; optional Java moss adds fine refuge. Keep food areas and intakes accessible. Review plant pest/product history before transfer. Separate plant entries and attachment methods are later work. \textbf{[CS-FARM; CS-COOP]}}
\Panel{Two tanks need separate decisions}{The 74--76 F discussion band intersects Co-Op's shrimp preference and trade Kuhli guidance. It is not a heater instruction; shared pH/GH/KH suitability remains unresolved. Match safe replacement-water practices and avoid abrupt pH/mineral corrections. No stocking count is established. \textbf{[CS-COOP; CS-KUH; CS-VET; proposed]}}
\Panel{Water record / test each tank and replacement supply}{Date / tank: \hrulefill\quad Temperature: \hrulefill\par
pH: \hrulefill\quad GH (dGH): \hrulefill\quad KH (dKH): \hrulefill\par
Ammonia / nitrite / nitrate (units): \hrulefill\par
Cycle evidence / water preparation / follow-up: \hrulefill}
\Sources{Sources (clickable): \href{https://www.aquariumcoop.com/blogs/aquarium/cherry-shrimp-care}{CS-COOP: Co-Op cherry care}; \href{https://www.theshrimpfarm.com/posts/neocaridina-shrimp-care-breeding/}{CS-FARM: Shrimp Farm care}; \href{https://www.theshrimpfarm.com/posts/understanding-general-hardness/}{CS-GH: Shrimp Farm GH}; \href{https://www.theshrimpfarm.com/posts/understanding-carbonate-hardness/}{CS-KH: Shrimp Farm KH}; \href{https://www.msdvetmanual.com/exotic-and-laboratory-animals/aquarium-fish/management-of-aquarium-fish}{CS-VET: MSD veterinary care}; \href{https://www.aquariumcoop.com/blogs/aquarium/care-guide-for-kuhli-loaches}{CS-KUH: Co-Op Kuhli care}. Scopes, derivations and unknowns: adjacent YAML worksheets.}{cherry-shrimp / tank-setup}
